\section{The Final Days and the Fourth Sunday: December 17 to 24}\label{sec:dec17-24}

On the morning of December~17, Advent turns a decisive corner. The calendar drops week numbers and counts down by date across eight privileged weekdays---December~17 through the morning of December~24---aimed directly at the Nativity of the Lord (\textit{Universal Norms on the Liturgical Year and the Calendar}, no.~42). These days outrank every optional or obligatory memorial (Table of Liturgical Days, no.~9; if an optional memorial occurs, the priest may take only its Collect while keeping the Advent readings and remaining prayers, GIRM, no.~355a). Preface~II of Advent replaces Preface~I at Mass, the weekday Gospels walk continuously through the infancy chapters of Matthew and Luke, and each evening from December~17 through December~23 culminates at the \textit{Magnificat} with one of the seven Great O~Antiphons.

\subsection{The Eight Privileged Weekdays: December 17 to 24}

Reading the weekday Mass Scriptures of December~17--24 in order (\textit{Ordo Lectionum Missae}, nos.~193--200) is like walking slowly through an illuminated gallery toward a lighted stable door. Each morning sets an Old Testament promise directly beside its Gospel fulfillment:

\begin{itemize}
  \item \textbf{December~17} (\textit{Genesis~49:2, 8--10; Psalm~72 [71]; Matthew~1:1--17}): Jacob blesses Judah with the promise that ``the sceptre shall not be taken away from Juda \dots\ till he come that is to be sent, and he shall be the expectation of nations.'' Matthew answers with the entire human genealogy of Jesus Christ, son of David, son of Abraham---naming kings and exiles, saints and sinners, Tamar, Rahab, Ruth, and the wife of Uriah, down to ``Joseph the husband of Mary, of whom was born Jesus, who is called Christ.'' God did not edit human frailty out of His family tree.
  \item \textbf{December~18} (\textit{Jeremiah~23:5--8; Psalm~72 [71]; Matthew~1:18--25}): Jeremiah announces the righteous Branch raised up for David whose name is ``The Lord our just one.'' Matthew narrates St.~Joseph's silent anguish and dream: ``Joseph, son of David, fear not to take unto thee Mary thy wife, for that which is conceived in her, is of the Holy Ghost.'' Joseph speaks not a single recorded word in the Gospels; his eloquence is immediate obedience.
  \item \textbf{December~19} (\textit{Judges~13:2--7, 24--25a; Psalm~71 [70]; Luke~1:5--25}): An angel announces the unexpected birth of Samson to Manoah's barren wife; Gabriel then stands beside the altar of incense to tell the aged priest Zechariah that Elizabeth will bear John, who will go before the Lord in the spirit and power of Elijah. Zechariah's doubting silence prepares his lips for praise.
  \item \textbf{December~20} (\textit{Isaiah~7:10--14; Psalm~24 [23]; Luke~1:26--38}): King Ahaz refuses to ask God for a sign, so the Lord gives His own: ``Behold a virgin shall conceive, and bear a son and his name shall be called Emmanuel.'' Luke then brings Gabriel to Nazareth, where Mary asks in pure wonder how this shall be and answers, ``Behold the handmaid of the Lord: be it done to me according to thy word.''
  \item \textbf{December~21} (\textit{Song of Songs~2:8--14 or Zephaniah~3:14--18a; Psalm~33 [32]; Luke~1:39--45}): The Bride hears the Beloved ``leaping upon the mountains, skipping over the hills,'' and Luke shows Mary rising in haste to cross the Judean hill country to Elizabeth's door. At the sound of Mary's voice, John leaps in the womb and Elizabeth cries out, ``Blessed art thou that hast believed.''
  \item \textbf{December~22} (\textit{1~Samuel~1:24--28; Canticle 1~Samuel~2:1, 4--8; Luke~1:46--56}): Hannah brings the young Samuel to the house of the Lord in Shiloh and sings, ``My heart hath rejoiced in the Lord.'' Beside her stands Mary singing the \textit{Magnificat}, praising the Mighty One who puts down the mighty from their seat, exalts the humble, and fills the hungry with good things.
  \item \textbf{December~23} (\textit{Malachi~3:1--4, 23--24 [Vulgate~3:1--4; 4:5--6]; Psalm~25 [24]; Luke~1:57--66}): Malachi foretells the messenger sent before the Lord's face and Elijah returning to turn the hearts of fathers to their children. Luke recounts the birth of John the Baptist: Zechariah writes ``John is his name,'' his tongue is loosed, and the hill country asks in awe, ``What an one, think ye, shall this child be?''
  \item \textbf{December~24, Morning} (\textit{2~Samuel~7:1--5, 8b--12, 14a, 16; Psalm~89 [88]; Luke~1:67--79}): On the final morning of Advent, God refuses David's offer to build a cedar temple and promises instead to build David an everlasting house. Zechariah, filled with the Holy Spirit, sings the \textit{Benedictus}, blessing the God of Israel who raises up a horn of salvation in David's house and visits us as ``the Orient from on high'' to guide our feet into the way of peace.
\end{itemize}

\subsection{The Seven Great O~Antiphons (December 17 to 23)}

At Evening Prayer each night from December~17 through December~23, the Roman Rite frames the \textit{Magnificat} with the seven Great Antiphons---called the \textit{O~Antiphons} because every one opens with a cry of wonder (\textit{O}) addressing Christ by a biblical title, recalls a saving deed of Israel's history, and closes with a single urgent imperative: \textit{Veni} (``Come''). The same seven invocations also furnish the Gospel Acclamation verses at weekday Mass across these seven days.

Taken in reverse order from December~23 back to December~17, the initial letters of their Latin titles---\textbf{E}mmanuel, \textbf{R}ex, \textbf{O}riens, \textbf{C}lavis, \textbf{R}adix, \textbf{A}donai, \textbf{S}apientia---spell the Latin words \textbf{ERO CRAS}: ``Tomorrow I will be [there].'' On December~23, the Church finishes her sevenfold plea and hears her Lord's hidden answer across the whole week.

\begin{threestable}{Date \& Title}{Scriptural Roots}{Plea of the Church}
Dec~17: \textit{O Sapientia} (O~Wisdom) & Sirach~24:3; Wisdom~8:1 & Ordering all things mightily and sweetly: come to teach us the way of prudence. \\
\addlinespace
Dec~18: \textit{O Adonai} (O~Lord and Ruler) & Exodus~3:2; 19:16--20; 6:6 & Appearing to Moses in the burning bush and giving the Law on Sinai: come to redeem us with an outstretched arm. \\
\addlinespace
Dec~19: \textit{O Radix Jesse} (O~Root of Jesse) & Isaiah~11:1, 10; 52:15 & Standing as an ensign of the peoples before whom kings keep silence: come to deliver us, delay no longer. \\
\addlinespace
Dec~20: \textit{O Clavis David} (O~Key of David) & Isaiah~22:22; 42:7; Apocalypse~3:7 & Opening where no one shuts and shutting where no one opens: come and lead the captive from the house of prison. \\
\addlinespace
Dec~21: \textit{O Oriens} (O~Dayspring / Rising Sun) & Malachi~4:2; Wisdom~7:26; Luke~1:78--79 & Brightness of eternal light and Sun of Justice: come and enlighten those who sit in darkness and the shadow of death. \\
\addlinespace
Dec~22: \textit{O Rex Gentium} (O~King of the Nations) & Haggai~2:7; Ephesians~2:14, 20; Genesis~2:7 & Desired of the nations and Cornerstone making two into one: come and save man whom Thou hast formed from the dust. \\
\addlinespace
Dec~23: \textit{O Emmanuel} (O~God-with-us) & Isaiah~7:14; 33:22; Genesis~49:10 & Our King and Lawgiver, expectation and Savior of the nations: come to save us, O~Lord our God. \\
\end{threestable}

The complete verified Latin texts (from the 1888 Ratisbon Roman Breviary) and classic English translations (by the Marquess of Bute, 1908) are printed in full in Appendix~B for nightly recitation at the Advent wreath.

\subsection{The Fourth Sunday of Advent Across Cycles A, B, and C}

Because the Fourth Sunday of Advent always falls between \textbf{December~18 and December~24}, it stands squarely inside the second phase of Advent. On that Sunday, the Fourth Sunday Mass replaces the date-keyed weekday readings for that day (and in the fourteen years per century when December~24 itself is Sunday---such as 2028 and 2034---the Fourth Sunday is celebrated on Sunday morning while Christmas Eve Evening Prayer~I and the Christmas Vigil Mass begin that evening). Across all three years, the Fourth Sunday sets Mary and Joseph before the altar on the very eve of Bethlehem:

\begin{threestable}{Cycle}{Appointed Readings}{Thematic Arc}
Year A & Isaiah~7:10--14; Psalm~24 (23); Romans~1:1--7; Matthew~1:18--24 & The Virgin bearing Emmanuel; the King of glory entering the gates; Jesus Christ son of David and Son of God; St.~Joseph's righteous obedience. \\
\addlinespace
Year B & 2~Samuel~7:1--5, 8b--12, 14a, 16; Psalm~89 (88); Romans~16:25--27; Luke~1:26--38 & God building David an eternal house; the mystery hidden for ages now made manifest; Gabriel's Annunciation and Mary's \textit{Fiat}. \\
\addlinespace
Year C & Micah~5:1--4a (Vulgate~5:2--5a); Psalm~80 (79); Hebrews~10:5--10; Luke~1:39--45 & Little Bethlehem yielding the eternal Ruler; ``a body thou hast fitted to me---Behold, I come to do thy will''; Mary hastening to Elizabeth at the Visitation. \\
\end{threestable}

\subsection{St.~Bernard, St.~Ambrose, and St.~Alphonsus on the Final Watch}

Three great doctors of the Church give voice to the quiet wonder of these days. Meditating on the Annunciation in his fourth \textit{Homily on the Missus est} (IV.8), St.~Bernard of Clairvaux imagines all creation---Adam in exile, Abraham, David, the angels, and every grieving sinner---holding its breath before the young girl of Nazareth:

\begin{quote}
The angel awaits your reply, for it is time that he should return to God, Who sent him. We, too, are waiting, O Lady, for a word of mercy---we, who are groaning under the sentence of condemnation. \dots\ Answer the word, receive the Word. Utter yours, conceive the Divine. Speak the word that is transitory, and embrace the Word that is everlasting. \dots\ Arise---run---open! Arise by faith, run by devotion, open by acceptance.
\end{quote}

In his second homily on the same Gospel (\textit{Missus est} II.7, II.9), Bernard compares Christ's quiet conception to dew falling silently upon Gideon's fleece (Psalm~71:6 [Hebrew~72:6]), marveling that in Mary's womb ``we behold light withholding its rays, the Word an infant, the Living Water athirst, Him Who is the Bread of Heaven suffering hunger.''

Preaching on the Visitation (Luke~1:39--45, proclaimed on December~21 and the Fourth Sunday of Year~C), St.~Ambrose of Milan (\textit{Exposition on Luke} II.19, II.26) draws out both the speed of charity and the universal motherhood of faith. Mary went up into the hill country neither doubting the angel nor testing Elizabeth: ``God was already in her womb, and her feeling bore her continually upward. The grace of the Holy Spirit knoweth no slow working'' (\textit{nescit tarda molimina sancti spiritus gratia}). Then Ambrose turns Mary's faith directly toward every believer (in the 1843 Oxford \textit{Catena Aurea} rendering): ``But let every one have the spirit of Mary, so that he may rejoice in the Lord. If according to the flesh there is one mother of Christ, yet, according to faith, Christ is the fruit of all.''

For daily meditation from December~16 through December~24, St.~Alphonsus Liguori composed his beloved \textit{Christmas Novena} (written in 1750 and published in 1758 in \textit{The Incarnation, Birth, and Infancy of Jesus Christ}), guiding ordinary Christians across nine short steps from the Father's gift of His Son and the hidden poverty of Jesus in Mary's womb to Joseph and Mary's road toward Bethlehem. Prayed beside the French Advent chant \textit{Rorate caeli} (``Drop down dew, ye heavens, from above,'' printed in Latin and English in Appendix~B) and the nightly O~Antiphon, these last days hush the soul so that when the bells of Christmas Night ring out, the heart is already kneeling.
